\begin{table}[t]
\caption{\textbf{Structural comparison between RAMSES and representative systems.} RAMSES uniquely enables full coordination across VRAM, DRAM, and NVMe, while prior systems provide only partial support. Abbreviations: Part.\ (partial support), Arch.\ (architectural-level analysis only), N/A (no empirical evaluation duration reported), and TP/W (throughput per watt).}

\centering
\label{tab:related_positioning}
\setlength{\tabcolsep}{2.5pt}
\renewcommand{\arraystretch}{1.1}
\begin{tabular}{lccccc}
\toprule
\textbf{System} &
\textbf{Tier} &
\textbf{NUMA} &
\textbf{Regime} &
\textbf{Eval.} &
\textbf{Energy}\\
&
\textbf{Scope} &
\textbf{aware} &
\textbf{Model} &
\textbf{Dur.} &
\textbf{Metric} \\
\midrule
FlexGen      & GPU--CPU--disk& $\times$ & None & Short & $\times$ \\
SwapAdvisor  & GPU--CPU     & $\times$ & None & Short & $\times$ \\
NEO          & Partial      & $\times$ & None & Short & $\times$ \\
SpecOffload  & GPU--NVMe    & $\times$ & None & Short & $\times$ \\
vLLM         & GPU-only     & $\times$ & None & Short & Part.    \\
Ind.\ Edge/DT& Arch.        & Part.    & None & N/A   & $\times$ \\
\textbf{RAMSES}
              & \textbf{Full} & $\checkmark$ & \textbf{Analytical} & \textbf{24--72h} & \textbf{TP/W} \\
\bottomrule
\end{tabular}
\end{table}

\subsection{Positioning of RAMSES}
\label{sec:positioning}

Table~\ref{tab:related_positioning} lays out the structural differences
between RAMSES and a set of representative systems along five axes:
tier scope, NUMA awareness, the formality of regime modeling,
evaluation duration (Eval.~Dur.), and energy reporting (throughput-per-watt,
TP/W). ``Part.'' marks partial support; ``Arch.'' marks
architectural-level analysis only.
The general picture is that prior LLM-serving frameworks tune for
throughput or stretch memory via partial GPU--CPU coordination---but
none of them evaluates the same combination of intra-node NUMA placement, resource preallocation, and pressure-driven policy switching.

Several recent TII papers cover complementary parts of the stack
without touching the intra-node memory dynamics that RAMSES targets.
\textbf{Tao~\emph{et~al.}~\cite{tao2019digitaltwin}} model latency
at the digital-twin synchronization layer, but they do not engage
GPU VRAM--DRAM swap dynamics or single-node pressure-driven inference.
\textbf{Cao~\emph{et~al.}~\cite{cao2021edgecps}} survey cloud--edge
resource coordination and explicitly bracket out the intra-node memory
hierarchy; RAMSES supplies the missing memory-subsystem layer.
\textbf{Chen~\emph{et~al.}~\cite{chen2021drl}} use deep RL to allocate
mobile-edge compute at the task-scheduling level. Inside each node,
RAMSES is what manages VRAM/DRAM/NVMe once the task has been placed.
\textbf{Wu~\emph{et~al.}~\cite{wu2021dnninference}} study
accuracy--latency tradeoffs through model splitting, with each device
treated as a black box. RAMSES empirically reduces the per-device delay that such
split-inference schedulers implicitly assume.
\textbf{Lu~\emph{et~al.}~\cite{lu2021dtfl}} reduce communication
overhead in federated learning for digital twins; on each participating
node, RAMSES reduces local inference latency and VRAM pressure, which
shortens per-round computation.
\textbf{Chekired~\emph{et~al.}~\cite{chekired2018fogscheduling}}
propose hierarchical fog-layer scheduling for smart-factory IoT data
and treat inference nodes as latency constants---which is exactly
where RAMSES plugs in by giving a regime-aware way to derive those
constants from VRAM--NVMe behavior.
\textbf{Deng~\emph{et~al.}~\cite{deng2020edgerl}} use RL for dynamic
edge-resource allocation, modeling inference cost at the
request-dispatch level; RAMSES supplies measured per-node latency distributions that an RL policy could use as observations.

The pattern is consistent: prior TII work coordinates at the network,
task, or device level and treats inference nodes as opaque units with
fixed latency. RAMSES sits one layer below all of this, characterizing
regime-dependent behavior across VRAM--DRAM--NVMe inside a single
serving node and producing empirical latency distributions that upstream schedulers can use.

\subsection{Distinct Contributions of RAMSES}

What is distinctive in RAMSES is the combination of two things that
usually appear separately. First, hierarchical memory orchestration.
Second, an analytical regime characterization built on a latency
decomposition over compute, locality, transfer, and synchronization.
Going past conventional serving optimizations, we identify
observable pressure ratios that classify capacity and transfer stress.
The three internal modules play distinct roles---the GPU Process
Allocator handles context conflicts, the Task Memory Manager handles
NUMA-aware allocation, and the GPU Booster handles swap-aware
asynchronous control---but each is judged against the same regime
boundaries. Where systems-venue work tends to optimize for throughput
or for one specific layer, the focus here is single-node QoS and sustained-load behavior under a parameterized synthetic workload.
